\specialchapter{记号索引}

参考编号依次表示章节、节和小节（或习题）。

\(l_{E}\)（E 为集合），U.V，UV（U、V 为加法子群），\(a^{0}\)（a 为理想）：第一章的预备约定

E:F:I.2.10

\(\mathbf{A}[\mathbf{S}^{-1}], a / s\)（A 为环，S 为 \(\mathbf{A}\) 的子集，\(a \in \mathbf{A}\)，\(s\) 为 \(\mathbf{S}\) 中元素的乘积）：II.2.1 \(i_{\mathbf{A}}^{\mathbf{S}}\) ：II.2.1

\(\mathbf{S}^{-1}\mathbf{A},\mathbf{A}_{p}\)（S 为乘性子集，p 为素理想）：II.2.1

\(\mathbf{M}[\mathbf{S}^{-1}], m / s, i_{\mathbf{M}}^{\mathbf{S}}\)（M 为 A-模，S 为 \(\mathbf{A}\) 的子集，\(m \in \mathbf{M}\)，\(s\) 为 \(\mathbf{S}\) 中元素的乘积）：II.2.2

\(S^{-1}M, M_{p}\)（M 为 A-模，S 为 A 的乘性子集，p 为 A 的素理想）：II.2.2

\(\mathbf{S}^{-1}u, u_{p}\)（\(u\) 为 A-模同态）：II.2.2

r(a)（a 为理想）：II.2.6

\(\mathbf{V}(\mathbf{M}), \mathbf{V}(f) (\mathbf{M}\) 是环 \(\mathbf{A}, f \in \mathbf{A})\)：II.4.3

\(\operatorname {Spec}(\mathbf{A})\) : II.4.3

\(\mathbf{X}_f\)（\(f \in \mathbf{A}\)，\(\mathbf{X} = \operatorname{Spec}(\mathbf{A})\)）：II.4.3

\(\Im (\mathbf{Y})\)（Y 是 \(\operatorname {Spec}(\mathbf{A}))\) 的子集）：II.4.3

\(^{a}h\)（h 为环同态）：II.4.3

Supp(M)（M 为 A-模）：II.4.4

\(\mathbf{A}_f, \mathbf{M}_f, u_f\)（A 为环，M 为 A-模，\(u\) 为 A-同态，\(f \in \mathbf{A}\)）：II.5.1

\(\mathbf{rg}_{\mathfrak{p}}(\mathbf{P})\)（P 为投射模）：II.5.3

rg(P)（P 为投射模）：II.5.3

\(P(A)\)，cl(M)（A 为环，M 为秩 1 的投射 A-模）：II.5.4

\(\operatorname{det}(u), \chi_u\)（\(u\) 为秩 \(n\) 的投射模的自同态）：II.5. 习题 9

\(\mathbf{A}^{(d)}, \mathbf{M}^{(d,k)}, \mathbf{M}^{(d)}\)（A 为分次环，M 为分次 A-模）：III.1.3

\(A_{(p)}, M_{(p)}\)（A 为分次环，p 为 A 的分次素理想，M 为分次 A-模）：III.1.4

\(\mathbf{gr}_n(\mathbf{G}), \mathbf{gr}(\mathbf{G})\)（G 为滤过群）：III.2.3

gr(h)（h 为与滤过相容的同态）：III.2.4

\(\mathbf{Z}_n\)（\(n\) 为整数 \(>1\)）：III.2.12

\(\hat{\mathbf{Z}}\) : III.2.13

\(\mathrm{A}\{\mathrm{X}_1,\dots ,\mathrm{X}_p\}\)（A 为线性拓扑环）：III.4.2

\(f(b_{1},\ldots ,b_{p})\)（f 为限制形式幂级数）：III.4.2

\(\mathbf{f} \circ \mathbf{g}, M_{t}, M_{t}(\mathbf{X}), J_{t}, J_{t}(\mathbf{X}), \mathbf{X}, \mathbf{1}_{n}\)（\(\mathbf{f}, \mathbf{g}\) 为形式幂级数系，\(\mathbf{g}\) 不含常数项）：III.4.4

\(\mathbf{f}(\mathbf{x})\)（f 为形式幂级数系，\(\mathbf{x}\) 为拓扑幂零元系）：III.4.5

\(\mathfrak{m}^{\times n}\)（m 为理想）：III.4.5

\(\operatorname{Ass}_{\mathbf{A}}(\mathbf{M}), \operatorname{Ass}(\mathbf{M})\)（M 为 A-模）：IV.1.1

\(\mathbf{Ass}_f(\mathbf{M})\) ：IV.1. 习题 17

\(\mathbf{A}^{\mathcal{G}}\)（A 为代数，\(\mathcal{G}\) 为作用在 A 上的群）：V.1.9

\(\mathcal{G}^{\mathrm{Z}}(\mathfrak{p}^{\prime}), \mathcal{G}^{\mathrm{Z}}, \mathrm{A}^{\mathrm{Z}}(\mathfrak{p}^{\prime}), \mathrm{A}^{\mathrm{Z}}\)（G 为作用在环 \(\mathrm{A}'\) 上的群，\(\mathfrak{p}'\) 为 \(\mathrm{A}'\) 的素理想）：V.2.2

\(\mathcal{G}^{\mathrm{T}}(\mathfrak{p}')\)，\(\mathcal{G}^{\mathrm{T}}\)，\(\mathbf{A}^{\mathrm{T}}(\mathfrak{p}')\)，\(\mathbf{A}^{\mathrm{T}}\)（\(\mathcal{G}\) 为作用在环 \(\mathbf{A}'\) 上的群，\(\mathfrak{p}'\) 为 \(\mathbf{A}'\) 的素理想）：V.2.2

\(\mathbf{K}^{\mathrm{Z}}(\mathfrak{p}')\)，\(\mathbf{K}^{\mathrm{Z}}\)，\(\mathbf{K}^{\mathrm{T}}(\mathfrak{p}')\)，\(\mathbf{K}^{\mathrm{T}}\)（K 是整闭整环

A 的分式域，\(p'\) 是 K 的拟 Galois 扩张中 A 的整闭包的素理想）：V.2.3

\(\mathbf{Y}^{\mathbf{p}}\)（其中 \(\mathbf{p} = (p_1, \ldots, p_m)\)，\(p_i\) 为整数 \(\geqslant 0\)）：V.3.1

\(\mathfrak{m}(\mathbf{A}), \kappa(\mathbf{A}), \mathrm{U}(\mathbf{A})\)（A 为局部环）：VI

\(\tilde{\mathbf{K}},\infty :\mathbf{VI}.2.1\)

+∞: VI.3.1

\(\Gamma_{\mathbf{A}}, v_{\mathbf{A}}\) : VI.3.2

\(\mathfrak{a}(\mathbf{M})\)（M 为优势集）：VI.3.5

\(h(\mathbf{G})\)（G 为全序群）：VI.4.4

\(\mathcal{T}_v\)（\(v\) 为赋值）：VI.5.2

\(e(v^{\prime} / v), e(\mathbf{A}^{\prime} / \mathbf{A}), e(\mathbf{L} / \mathbf{K})\) : VI.8.1

\(f(v^{\prime} / v),f(\mathbf{A}^{\prime} / \mathbf{A}),f(\mathbf{L} / \mathbf{K})\) : VI.8.1

\(\varepsilon (\mathrm{G},\mathrm{H})\)（G 为全序群，H 为 G 的有限指数子群）：VI.8.4

\(\varepsilon(v'/v)\)（\(v\) 为赋值，\(v'\) 为 \(v\) 的扩张）：VI.8.4

mod(x), mod\_K(x)（K 为非离散局部紧域，x ∈ K）：VI.9.1

\(r(\mathbf{G})\)（交换群的有理秩）：VI.10.2

\(d(\mathbf{K}' / \mathbf{K}), s(v' / v), r(v' / v)\)（v 为 \(\mathbf{K}\) 上的赋值，\(v'\) 为 \(v\) 在超越扩张 \(\mathbf{K}'\)（\(\mathbf{K}\) 的）上的扩张）：VI.10.3

I(A), D(A)（A 为整环）：VII.1.1

\(\mathfrak{a} \prec \mathfrak{b}\)，\(\operatorname{div}(\mathfrak{a})\)，\(\operatorname{div}(x)\)（\(\mathfrak{a}\)、\(\mathfrak{b}\) 为分式理想，\(x\) 为分式域的元素）：VII.1.1

\(\tilde{\mathfrak{a}}\)（a 为分式理想）：VII.1.1

\(d_{1} \leqslant d_{2}\)（\(d_{1}, d_{2}\) 为除子）：VII.1.1

b: a（a、b 为分式理想）：VII.1.1

J(A)（A 为整环）：VII.1.2

P(A)（A 为 Krull 整环）：VII.1.3

\(\mathfrak{p}^{(n)}\)（p 为除性素理想）：VII.1.4

\(v_{p}\)（p 为 Krull 整环中高度 1 的素理想）：VII.1.10

\(\mathbf{F}(\mathbf{A}), \mathbf{C}(\mathbf{A})\)（A 为 Krull 整环）：VII.1.10

\(e(\mathfrak{P}/\mathfrak{p})\) \((\mathfrak{p}\in\mathrm{P}(\mathrm{A}),\;\mathfrak{P}\in\mathrm{P}(\mathrm{B}),\;\mathrm{A}\;\text{和}\;\mathrm{B}\;\text{Krull 整环},\;\mathrm{A}\subset\mathrm{B},\;\mathfrak{P}\cap\mathrm{A}=\mathfrak{p}):\) VII.1.10

\(i\)（从 \(\mathbf{D}(\mathbf{A})\) 到 \(\mathbf{D}(\mathbf{B})\) 的同态，或从 \(\mathbf{C}(\mathbf{A})\) 到 \(\mathbf{C}(\mathbf{B}))\) 的同态）：VII.1.10

\(\bar{i}\)（从 \(\mathbf{C}(\mathbf{A})\) 到 \(\mathbf{C}(\mathbf{B}))\) 的同态）：VII.1.10

\(A, A_{0}, \Delta(K)\)（限制阿代尔的环）：VII.2.4

\(\mathfrak{P}^*, \mathfrak{P}^*(\mathbf{A})\)（A 为整环）：VII.3.2

M*（格 M 的对偶格）：VII.4.2

\(l_{\mathfrak{p}}(\mathrm{T}), \chi (\mathrm{T})\)（T 为扭 A-模，\(\mathfrak{p}\) 为高度 1 的素理想）：VII.4.5

\(\pmb {F}(\mathbf{A}),\pmb{T}(\mathbf{A}),\mathrm{cl}(\mathbf{M}):\) VII.4.5

\(\chi (\mathbf{M},\mathbf{M}^{\prime})\)（M、\(\mathbf{M}'\) 为格）：VII.4.6

\(c(d)\)（\(d\) 为除子）：VII.4.7

\(c(\mathbf{M}), r(\mathbf{M}), \gamma(\mathbf{M})\)（M 为格）：VII.4.7

\(\mathfrak{P}|\mathfrak{p},e_{\mathfrak{P} / \mathfrak{p}},f_{\mathfrak{P} / \mathfrak{p}},f(\mathfrak{P} / \mathfrak{p})\)（A \(\subset\) B 为 Krull 整环，且 B 是有限 A-代数

\(\mathfrak{p} \in \mathrm{P}(\mathbf{A}), \mathfrak{P} \in \mathrm{P}(\mathbf{B}), \mathfrak{P} \cap \mathbf{A} = \mathfrak{p}\) : VII.4.8

\(\mathbf{N}_{\mathbf{B} / \mathbf{A}},\mathbf{N},i_{\mathbf{B} / \mathbf{A}}\) :VII.4.8

